\subsubsection{Simulated Annealing}

In the third difficulty, 
we use Simulated Annealing to plan the route for the monster AI.
Although Simulated Annealing can not guarantee finding the shortest path,
it can still be used to obtain a near-optimal path, 
which is sufficient for our Monster AI.

Initially, the initial state is built by greedy search
with toroidal Manhattan distance as its heuristic function.
However, after some tests, we found that in some cases,
greedy search might fail to reach the target,
causing the entire simulated annealing process to fail.
Therefore, we used Depth-First Search(DFS) to generate the initial state.
The DFS prefers moves that are closer to the human using toroidal Manhattan distance,
but it also shuffles candidates so the path is not completely deterministic.
Several tests proved that this metric is effective enough to ensure reaching the target state,
and the time complexity is also acceptable because of its preference for smaller Manhattan distance.

The Simulated Anealing generates neibour solution by cutting the current path at a random point 
and trying to rebuild the remaining path from that cut point to the human.
The cost of each state is calculated by the following function:
\[
S(P)=|P|+
\begin{cases}
0, & p_{\mathrm{end}} = h,\\
50D(p_{\mathrm{end}},h)+1000, & p_{\mathrm{end}} \neq h.
\end{cases}
\]
